% unidades/13-te-form-1.tex
\chapter{🔗 て形 (1) — La Forma て}
\unidadheader{13}{て形 (1)}{La Forma て}{Conectar acciones y pedir favores}

\begin{tcolorbox}[vocabbox, title={📌 Vocabulario Nuevo}]
\begin{tabularx}{\linewidth}{llX}
\toprule
\textbf{Japonés} & \textbf{Español} & \textbf{Tipo} \\
\midrule
\vocabitem{待つ}{まつ}{esperar}{v. gr.1}
\vocabitem{入る}{はいる}{entrar}{v. gr.1}
\vocabitem{出す}{だす}{sacar / entregar}{v. gr.1}
\vocabitem{持つ}{もつ}{tener / sostener}{v. gr.1}
\vocabitem{使う}{つかう}{usar}{v. gr.1}
\vocabitem{座る}{すわる}{sentarse}{v. gr.1}
\vocabitem{立つ}{たつ}{pararse}{v. gr.1}
\vocabitem{教える}{おしえる}{enseñar}{v. gr.2}
\vocabitem{見せる}{みせる}{mostrar}{v. gr.2}
\vocabitem{開ける}{あける}{abrir}{v. gr.2}
\bottomrule
\end{tabularx}
\end{tcolorbox}

\subsection*{🗣️ Frases Clave}

\frase{ちょっと\ruby{待}{ま}ってください。}{Espere un momento, por favor.}
\frase{ここに\ruby{座}{すわ}ってください。}{Siéntese aquí, por favor.}
\frase{窓を\ruby{開}{あ}けてもいいですか？}{¿Puedo abrir la ventana?}
\frase{はい、どうぞ。}{Sí, adelante.}

\subsection*{⚙️ Gramática}

\patron{Reglas de la forma て (Grupo 1)}
  {う・つ・る → って \quad | \quad ぬ・ぶ・む → んで \quad | \quad く → いて \quad | \quad ぐ → いで \quad | \quad す → して}
  {\ej{\ruby{待}{ま}つ → \ruby{待}{ま}って}
  \ej{\ruby{飲}{の}む → \ruby{飲}{の}んで}}

\patron{Pedir un favor: 〜てください}
  {V-て + ください}
  {\ej{\ruby{名前}{なまえ}を\ruby{書}{か}いてください。}{Escriba su nombre, por favor.}
  \ej{ちょっと\ruby{待}{ま}ってください。}{Espere un poco, por favor.}}

\begin{tcolorbox}[ejerciciobox, title={📝 Ejercicios Rápidos}]
\begin{enumerate}
  \item Convierte a forma て: \ruby{飲}{の}む.
  \item Convierte a forma て: \ruby{話}{はな}す.
  \item Di "Présteme esto, por favor" (\ruby{貸}{か}す).
\end{enumerate}
\vspace{0.4em}
\textbf{Respuestas:}
1) \ruby{飲}{の}んで\quad
2) \ruby{話}{はな}して\quad
3) これを\ruby{貸}{か}してください。
\end{tcolorbox}

\begin{tcolorbox}[kanjibox, title={🀄 Kanjis de la Unidad}]
\begin{tabularx}{\linewidth}{cllXX}
\toprule
\textbf{Kanji} & \textbf{On} & \textbf{Kun} & \textbf{Significado} & \textbf{Vocab} \\
\midrule
\kanjirow{入}{にゅう}{はい・い}{entrar}{\ruby{入}{はい}る / \ruby{入口}{いりぐち}}
\kanjirow{出}{しゅつ}{で・だ}{salir / sacar}{\ruby{出}{で}る / \ruby{出口}{でぐち}}
\kanjirow{持}{じ}{も}{sostener / tener}{\ruby{持}{も}つ / \ruby{持}{も}って\ruby{来}{く}る}
\kanjirow{借}{しゃく}{か}{pedir prestado}{\ruby{借}{か}りる / \ruby{借金}{しゃっきん}}
\kanjirow{洗}{せん}{あら}{lavar}{\ruby{洗}{あら}う / \ruby{洗濯}{せんたく}}
\bottomrule
\end{tabularx}
\end{tcolorbox}
